%% FullCourse_10Lecs -- Lecture 9, Topic 9.2: Under the Hood + First Session + 5-Pillar Workflow
%% Source: ./archive_pre_split/Lec09_Agentic_AI_v3.tex (sections 3-5)
%% Pass B: layer in refinements from
%%         ../../MiniCourse_8hr/Slides/Module4_Agentic_AI/M4_T4_Claude_Code_Operating_Model.tex
%% Rewritten 2026-07-06: Act 2 "Machine" of the five-act structure. Adds install/init
%%         in main flow, CLI anatomy, context-window + compaction mechanics; first
%%         session retargeted to the ES Fellows roster (Lec10 T4 thread). Tool
%%         mechanics fact-checked against code.claude.com docs (2026-07).
%%         See Lec09_Rewrite_Plan_2026-07-06.md.

%% Shared preamble for the FullCourse_10Lecs topic decks.
%%
%% Each topic .tex \input's this file, then defines its own \title, \date
%% override (if any), \graphicspath, and \begin{document}.
%%
%% Reconstructed verbatim from the inlined copy preserved in
%% Lec10_Case_Studies/Lec10_T8_Case6_DeepRamsey.tex (lines 23-190), which is
%% explicitly annotated there as "from ../shared_preamble.tex, copied verbatim".
%% Base preamble matches MiniCourse_8hr/Slides/shared_preamble.tex; this version
%% additionally provides \sectiondivider, the conceptbox/cautionbox/examplebox/
%% takeawaybox tcolorboxes, and \compacttables used across the split decks.

\documentclass[10pt,english,aspectratio=169]{beamer}

\usepackage{ae,aecompl}
\usepackage[T1]{fontenc}
\usepackage{textcomp}
\usepackage[utf8]{inputenc}
\usepackage{babel}
\usepackage{datetime2}

\setcounter{secnumdepth}{3}
\setcounter{tocdepth}{3}

\usepackage{amsmath, amssymb, amsthm, graphicx}
\usepackage{booktabs, multirow}
\usepackage{tikz}
\usepackage{pgfplots}
\pgfplotsset{compat=1.16}
\usepackage[most]{tcolorbox}
\tcbuselibrary{skins, breakable, theorems}
\usepackage{listings}
\usepackage{xcolor}

\lstset{
  basicstyle=\ttfamily\small,
  keywordstyle=\color{blue!80!black}\bfseries,
  commentstyle=\color{green!50!black}\itshape,
  stringstyle=\color{orange!80!black},
  backgroundcolor=\color{gray!8},
  frame=single,
  rulecolor=\color{gray!40},
  breaklines=true,
  showstringspaces=false,
  numbers=none,
}

\lstdefinestyle{pythonstyle}{
    language=Python,
    basicstyle=\ttfamily\footnotesize,
    keywordstyle=\color{blue!80!black}\bfseries,
    commentstyle=\color{green!50!black}\itshape,
    stringstyle=\color{orange!80!black},
    frame=single,
    breaklines=true,
    showstringspaces=false,
    numbers=left,
    numbersep=5pt,
    numberstyle=\tiny\color{gray},
    backgroundcolor=\color{gray!8},
    rulecolor=\color{gray!40}
}

\ifx\hypersetup\undefined
  \AtBeginDocument{%
    \hypersetup{bookmarks=false,breaklinks=true,pdfborder={0 0 0},colorlinks=true,
      linkcolor=DarkText,urlcolor=RedTitle}%
  }
\else
  \hypersetup{bookmarks=false,breaklinks=true,pdfborder={0 0 0},colorlinks=true,
    linkcolor=DarkText,urlcolor=RedTitle}%
\fi

\makeatletter
\newcommand\makebeamertitle{%
  \begin{frame}[plain]
    \vspace*{1.0cm}
    \centering
    {\usebeamerfont{title}\usebeamercolor[fg]{title}\inserttitle\par}
    \vspace{1.0cm}
    {\usebeamerfont{author}\usebeamercolor[fg]{author}\insertauthor\par}
    \vspace{0.2cm}
    {\usebeamerfont{date}\usebeamercolor[fg]{date}\insertdate\par}
  \end{frame}}%
\AtBeginDocument{%
  \let\origtableofcontents=\tableofcontents
  \def\tableofcontents{\@ifnextchar[{\origtableofcontents}{\gobbletableofcontents}}
  \def\gobbletableofcontents#1{\origtableofcontents}%
}

\usetheme{metropolis}
\usefonttheme{professionalfonts}

\definecolor{LightGray}{RGB}{230,230,230}
\definecolor{RedTitle}{RGB}{0,0,200}
\definecolor{VeryLightGray}{RGB}{245,245,245}
\definecolor{DarkText}{RGB}{0,0,0}

\setbeamercolor{background canvas}{bg=white}
\setbeamercolor{title}{fg=RedTitle}
\setbeamercolor{author}{fg=DarkText}
\setbeamercolor{date}{fg=DarkText}
\setbeamercolor{frametitle}{bg=VeryLightGray,fg=DarkText}
\setbeamercolor{normal text}{fg=DarkText}
\setbeamerfont{title}{size=\large}

\setbeamersize{text margin left=5mm, text margin right=5mm}
\addtobeamertemplate{frametitle}{}{\vspace{-0.5em}}
\newcommand{\reducevspace}{\vspace{-0.5cm}}

% Shared section divider and box styles for split topic decks.
\newcommand{\sectiondivider}[2]{%
\begin{frame}[plain,noframenumbering]
\begin{center}
\vfill
{\Large\textbf{#1}\par}
\vspace{0.35cm}
{\normalsize #2\par}
\vfill
\end{center}
\end{frame}
}

\newtcolorbox{conceptbox}[1][]{
  colback=blue!3,
  colframe=RedTitle!55,
  boxrule=0.35mm,
  arc=1mm,
  left=2mm,
  right=2mm,
  top=1mm,
  bottom=1mm,
  #1
}
\newtcolorbox{cautionbox}[1][]{
  colback=red!3,
  colframe=red!50!black,
  boxrule=0.35mm,
  arc=1mm,
  left=2mm,
  right=2mm,
  top=1mm,
  bottom=1mm,
  #1
}
\newtcolorbox{examplebox}[1][]{
  colback=green!3,
  colframe=green!45!black,
  boxrule=0.35mm,
  arc=1mm,
  left=2mm,
  right=2mm,
  top=1mm,
  bottom=1mm,
  #1
}
\newtcolorbox{takeawaybox}[1][]{
  colback=VeryLightGray,
  colframe=RedTitle!55,
  boxrule=0.35mm,
  arc=1mm,
  left=2mm,
  right=2mm,
  top=1mm,
  bottom=1mm,
  #1
}
\newcommand{\compacttables}{\renewcommand{\arraystretch}{1.15}}

\setbeamertemplate{navigation symbols}{}
\setbeamertemplate{footline}{%
  \hfill%
  \usebeamercolor[fg]{page number in head/foot}%
  \usebeamerfont{page number in head/foot}%
  \insertframenumber\,/\,\inserttotalframenumber\kern1em\vskip2pt%
}

\makeatother

\author{Zhigang Feng}
% "date with time": compile timestamp so each build is identifiable.
\date{\today\ \DTMcurrenttime}


% Suppress redundant auto-generated metropolis section page; \sectiondivider frames are the dividers we keep.
\metroset{sectionpage=none}

\graphicspath{{./pic/}{./figures/}{../../MiniCourse_8hr/Slides/Module4_Agentic_AI/pic/}{../}}


\usetikzlibrary{positioning,arrows.meta,shapes,calc}

%% Lec09 shared MAP slide: "one map, five acts" recurring diagram.
%% Input this file AFTER ../shared_preamble.tex (needs RedTitle, VeryLightGray)
%% and call \LecNineMap{k} inside a frame, where k in {1..5} highlights the
%% current act ("you are here") and k = 0 renders the closing reprise with all
%% five acts checked (used by T5's "Your Job Description" frame).
%% Filename deliberately does not match the build glob Lec*_T*.tex, so the
%% build script never tries to compile it standalone.

\newcommand{\LecNineMap}[1]{%
% Precompute per-box style selectors and the checkmark token; TikZ option
% lists are comma-split before TeX conditionals run, so \ifnum must not sit
% inside a [...] options list.
\ifnum#1=0
  \tikzset{lecmapa/.style={lecmapdone}, lecmapb/.style={lecmapdone},
           lecmapc/.style={lecmapdone}, lecmapd/.style={lecmapdone},
           lecmape/.style={lecmapdone}}%
  % green fill + the "all five acts complete" banner mark completion; a per-box
  % checkmark overflows the MANAGEMENT box, so keep the token empty here too.
  \def\lecmapchk{}%
\else
  \tikzset{lecmapa/.style={}, lecmapb/.style={}, lecmapc/.style={},
           lecmapd/.style={}, lecmape/.style={}}%
  \def\lecmapchk{}%
  \ifnum#1=1 \tikzset{lecmapa/.style={lecmapcur}}\fi
  \ifnum#1=2 \tikzset{lecmapb/.style={lecmapcur}}\fi
  \ifnum#1=3 \tikzset{lecmapc/.style={lecmapcur}}\fi
  \ifnum#1=4 \tikzset{lecmapd/.style={lecmapcur}}\fi
  \ifnum#1=5 \tikzset{lecmape/.style={lecmapcur}}\fi
\fi
\begin{center}
\begin{tikzpicture}[
    act/.style={rectangle, rounded corners, draw=black!60, fill=VeryLightGray,
        minimum width=2.6cm, minimum height=0.92cm, text width=2.5cm,
        align=center, font=\scriptsize, text=DarkText},
    lecmapcur/.style={draw=RedTitle, very thick, fill=orange!20},
    lecmapdone/.style={draw=green!45!black, thick, fill=green!10},
    q/.style={font=\tiny\itshape, text width=2.7cm, align=center, text=black!70},
    sat/.style={rectangle, rounded corners, draw=black!45, dashed, fill=white,
        text width=5.0cm, align=center, font=\tiny, text=black!70,
        minimum height=0.5cm},
    barlab/.style={font=\tiny\bfseries, text=black!60},
    arr/.style={-{Stealth[length=2mm]}, thick, gray!60}
]
    % --- five act boxes ---
    \node[act, lecmapa] (a1) at (0,0)
        {\textbf{9.1 MAP}\lecmapchk\\ \tiny what \& why};
    \node[act, lecmapb] (a2) at (2.85,0)
        {\textbf{9.2 MACHINE}\lecmapchk\\ \tiny under the hood};
    \node[act, lecmapc] (a3) at (5.7,0)
        {\textbf{9.3 METHOD}\lecmapchk\\ \tiny homotopy workflow};
    \node[act, lecmapd] (a4) at (8.55,0)
        {\textbf{9.4 MANAGEMENT}\lecmapchk\\ \tiny run the project};
    \node[act, lecmape] (a5) at (11.4,0)
        {\textbf{9.5 MINDSET}\lecmapchk\\ \tiny your role};
    \draw[arr] (a1) -- (a2);
    \draw[arr] (a2) -- (a3);
    \draw[arr] (a3) -- (a4);
    \draw[arr] (a4) -- (a5);
    % --- the student question each act answers ---
    \node[q] at (0,-0.82)    {what is this,\\ and why care?};
    \node[q] at (2.85,-0.82) {how does it work?\\ how do I start?};
    \node[q] at (5.7,-0.82)  {how do I structure\\ real work?};
    \node[q] at (8.55,-0.82) {how do I run a\\ whole project?};
    \node[q] at (11.4,-0.82) {what goes wrong?\\ what is my job?};
    % --- "you are here" marker above the current act ---
    \ifnum#1=1 \node[font=\scriptsize\bfseries, text=RedTitle] at (0,0.95)
        {you are here\ $\blacktriangledown$};\fi
    \ifnum#1=2 \node[font=\scriptsize\bfseries, text=RedTitle] at (2.85,0.95)
        {you are here\ $\blacktriangledown$};\fi
    \ifnum#1=3 \node[font=\scriptsize\bfseries, text=RedTitle] at (5.7,0.95)
        {you are here\ $\blacktriangledown$};\fi
    \ifnum#1=4 \node[font=\scriptsize\bfseries, text=RedTitle] at (8.55,0.95)
        {you are here\ $\blacktriangledown$};\fi
    \ifnum#1=5 \node[font=\scriptsize\bfseries, text=RedTitle] at (11.4,0.95)
        {you are here\ $\blacktriangledown$};\fi
    \ifnum#1=0 \node[font=\scriptsize\bfseries, text=green!45!black] at (5.7,0.95)
        {all five acts complete};\fi
    % --- bottom band: the three questions of the lecture ---
    \fill[blue!10]   (-1.3,-1.55) rectangle (4.15,-1.22);
    \fill[green!10]  (4.4,-1.55)  rectangle (9.85,-1.22);
    \fill[orange!15] (10.1,-1.55)  rectangle (12.7,-1.22);
    \node[barlab] at (1.425,-1.385)  {HOW IT WORKS};
    \node[barlab] at (7.125,-1.385)  {HOW TO USE IT WELL};
    \node[barlab] at (11.4,-1.385)   {YOUR ROLE};
    % --- dashed satellites: the two decks outside these five acts ---
    \node[sat] at (2.0,-2.2)
        {\textit{after 9.1:} optional interlude \textbf{9.1b} --- the agent landscape (MCP, A2A, benchmarks, safety)};
    \node[sat] at (9.8,-2.2)
        {\textit{after 9.5:} \textbf{9.6} wrap-up $\to$ \textbf{Lec10} full case studies (same morning)};
\end{tikzpicture}
\end{center}
}


\title[AI for Econ Research]{AI for Economic Research: Dynamic Models, Language, and Agents\\
Lecture 9.2: Under the Hood --- How It Works and How to Start}

\begin{document}

\makebeamertitle

%=============================================================================
\begin{frame}{The Map: Act 2 --- The Machine}
\textbf{This deck answers the second question: how does it actually work, and how do I start?}

\LecNineMap{2}

\vspace{0.1cm}
\footnotesize\textit{Route: the three physical components and the tool-call loop $\to$ install and initialize $\to$ the shape of the CLI $\to$ your first session (on the fellows roster) $\to$ tokens, the context window, and the pillars that keep you honest.}
\end{frame}

%=============================================================================
\section{The Machine}
%===========================================================

\sectiondivider{The Machine}{Harness, model, files --- and the loop that joins them}

\begin{frame}{Three Components: Harness, LLM, Your Files}

A terminal agent has three physically distinct parts --- the same anatomy in Claude Code, Codex CLI, Gemini CLI, or Aider:

\vspace{0.2cm}

\begin{columns}[T]
\begin{column}{0.55\textwidth}
\begin{enumerate}
    \item \textbf{The Harness} (lives on your machine)\\
    \small CLI software you install. Reads your files, runs shell commands, manages the conversation loop.

    \vspace{0.2cm}

    \item \textbf{The LLM} (lives on the provider's servers)\\
    \small Every message is bundled with conversation history and sent via HTTPS to the API. The ``thinking'' happens in the cloud.

    \vspace{0.2cm}

    \item \textbf{Your Files} (live on your machine)\\
    \small Python scripts, data CSVs, LaTeX---all on your disk. Code executes on your hardware with your packages.
\end{enumerate}
\end{column}
\begin{column}{0.4\textwidth}
\begin{center}
\begin{tikzpicture}[
    macbox/.style={rectangle, draw=blue!60, fill=blue!8, rounded corners,
                   minimum width=3.2cm, minimum height=1.0cm,
                   text width=3.0cm, align=center, font=\small},
    cloudbox/.style={rectangle, draw=orange!70, fill=orange!10, rounded corners,
                     minimum width=3.2cm, minimum height=1.0cm,
                     text width=3.0cm, align=center, font=\small},
    arrow/.style={-{Stealth[length=2mm]}, thick, gray}
]
    \node[macbox] (harness) {Harness + Files\\{\footnotesize Your machine}};
    \node[cloudbox, below=1.4cm of harness] (llm) {LLM\\{\footnotesize Anthropic API}};
    \draw[arrow, bend left=15] (harness) to node[right,font=\tiny] {prompt + context} (llm);
    \draw[arrow, bend left=15] (llm) to node[left,font=\tiny] {tool calls + text} (harness);
\end{tikzpicture}
\end{center}
\vspace{0.15cm}
\footnotesize
\textbf{Key:} code can run locally and many project files stay on your machine. That does \textit{not} mean the whole workflow is automatically private.
\end{column}
\end{columns}
\end{frame}

\begin{frame}[fragile]{The Loop in Machine Language: JSON Tool Calls}
\textbf{Lec07 T3a explained that an LLM can emit structured JSON via schema-constrained decoding. In a terminal agent, every tool call uses that mechanism.}

\textbf{Example.} When the agent decides to read a file, it emits:
\begin{lstlisting}[style=pythonstyle,language={},basicstyle=\ttfamily\scriptsize,numbers=none]
{"tool": "Read", "args": {"file_path": "fellows_seed.csv"}}
\end{lstlisting}

The harness parses the JSON, executes \texttt{Read}, and feeds the file content back as the next message; the agent reads it and decides the next action.

\begin{itemize}
    \item \textbf{Loop.} Model output (JSON) $\rightarrow$ harness execution $\rightarrow$ tool result $\rightarrow$ next model turn. Continues until the model emits a ``done'' signal or hits a stop condition.
    \item \textbf{Schema is the contract.} Each tool advertises an argument schema; outputs that fail the schema are rejected.
    \item \textbf{Same brain.} There is no separate ``planner'' --- the same transformer that writes code also decides to call tools.
    \item \textbf{Key implication.} Every action the agent takes is a \emph{next-token prediction} constrained by the available tool schema.
\end{itemize}
\end{frame}

\begin{frame}{Privacy Boundary: Local Execution Is Not the Same as Local Processing}
\begin{itemize}
    \item code may execute locally on your machine
    \item many project files may stay on disk locally
    \item if file contents, prompts, or tool outputs are sent back in the conversation, those contents may still traverse the network
    \item only a fully local or institutionally private deployment changes that boundary
\end{itemize}

\vspace{0.25cm}
\textbf{Practical rule:} running code locally does not make it safe to put restricted data into the conversation --- treat everything you type or attach as leaving your machine.
\end{frame}

% MiniCourse refinement: terminal-workflow vocabulary from M4_T4
\begin{frame}{Terminal Vocabulary in 60 Seconds}
\small
\begin{center}
\begin{tabular}{|p{2.7cm}|p{8.5cm}|}
\hline
\textbf{Term} & \textbf{Meaning} \\ \hline
\textbf{Terminal} & The text window where you type commands instead of clicking buttons. \\ \hline
\textbf{Shell} & The program that interprets commands. macOS often uses \texttt{zsh}; many examples are written in \texttt{bash}. \\ \hline
\textbf{Bash} & A common shell language. When Claude Code says it will use \texttt{Bash}, it means ``run this command in the terminal.'' \\ \hline
\textbf{Working directory} & The folder commands run from. \texttt{cd my-project} changes it; \texttt{pwd} prints it. \\ \hline
\textbf{PATH} & The list of folders your shell searches when you type a command such as \texttt{python} or \texttt{claude}. \\ \hline
\end{tabular}
\end{center}

\vspace{0.15cm}
\textbf{Why this matters:} terminal agents act through shell commands. If the working directory or \texttt{PATH} is wrong, the agent may be competent but operating in the wrong place.
\end{frame}

\begin{frame}{How the Terminal Pieces Fit Together}
\textbf{One picture for the five words: you type into the terminal; the shell finds the program via PATH and runs it in the working directory.}

\vspace{0.15cm}

\begin{center}
\begin{tikzpicture}[
    tb/.style={rectangle, rounded corners, draw=black!60, fill=VeryLightGray,
        minimum width=2.6cm, minimum height=1.05cm, text width=2.5cm,
        align=center, font=\scriptsize},
    side/.style={rectangle, rounded corners, draw=RedTitle!60, fill=blue!5,
        minimum width=3.4cm, minimum height=0.85cm, text width=3.3cm,
        align=center, font=\scriptsize},
    arr/.style={-{Stealth[length=2mm]}, thick, gray!60},
    lab/.style={font=\tiny\ttfamily, text=black!60}
]
    \node[tb] (you) at (0,0) {\textbf{You}\\ type a command:\\ \texttt{claude}};
    \node[tb] (term) at (3.3,0) {\textbf{Terminal}\\ the window that carries text in and out};
    \node[tb] (shell) at (6.6,0) {\textbf{Shell} (\texttt{zsh}/\texttt{bash})\\ interprets the line};
    \node[tb] (prog) at (9.9,0) {\textbf{Program runs}\\ \texttt{claude}, \texttt{python}, \texttt{git}};
    \node[side] (path) at (6.6,1.35) {\textbf{PATH:} the list of folders searched to \emph{find} the program};
    \node[side] (wd) at (9.9,-1.35) {\textbf{Working directory:} the folder the program runs \emph{in} (reads/writes here)};
    \draw[arr] (you) -- (term);
    \draw[arr] (term) -- (shell);
    \draw[arr] (shell) -- node[lab, above] {found it} (prog);
    \draw[arr, dashed] (path) -- (shell);
    \draw[arr, dashed] (wd) -- (prog);
\end{tikzpicture}
\end{center}

\begin{itemize}
    \item When you type \texttt{claude}, the shell finds it via \texttt{PATH} and starts it \emph{in your current folder} --- that is why \texttt{cd my-project} comes first
    \item The agent's \texttt{Bash} tool is this same picture one level down: the agent typing into the same shell, in the same working directory
\end{itemize}
\end{frame}

%===========================================================
%=============================================================================
\section{Start Here}
%===========================================================

\sectiondivider{Start Here}{Install, initialize, and the shape of the CLI}

\begin{frame}[fragile]{Install Once: Three Durable Steps}
\textbf{Setup has exactly three steps, and only the first command ever changes.}

\vspace{0.2cm}

\begin{enumerate}
    \item \textbf{Install the CLI} --- one installer command for your OS (native installer script; Homebrew and \texttt{npm} also work):
\begin{lstlisting}[style=pythonstyle,language={},basicstyle=\ttfamily\scriptsize,numbers=none]
curl -fsSL https://claude.ai/install.sh | bash
\end{lstlisting}
    \item \textbf{Authenticate once} --- run \texttt{claude}; a browser window opens; log in with your account or API key
    \item \textbf{Work from the project} --- every session starts the same way:
\begin{lstlisting}[style=pythonstyle,language={},basicstyle=\ttfamily\scriptsize,numbers=none]
cd my-project
claude
\end{lstlisting}
\end{enumerate}

\vspace{0.15cm}

\begin{takeawaybox}
\footnotesize
\textbf{Install once per machine. Authenticate once per account. Initialize every session per project --- automatically (next frames).}
\end{takeawaybox}

\vspace{0.1cm}
\footnotesize\textit{Durable vs.\ dated: this 3-step shape survives every product update; exact commands, subscription tiers, and the Windows/OneDrive caveats are dated detail $\to$ T6 Appendix A and the official docs.}
\end{frame}

\begin{frame}{Files First: Global vs.\ Local}

Two layers of configuration. \textbf{Local overrides global.}

\vspace{0.2cm}

\begin{center}
\footnotesize
\begin{tabular}{|l|l|p{4.9cm}|}
\hline
\textbf{Path (Claude Code example)} & \textbf{Scope} & \textbf{What lives here} \\ \hline
\texttt{\textasciitilde/.claude/settings.json} & Global & Default permissions, model preference \\ \hline
\texttt{\textasciitilde/.claude/CLAUDE.md} & Global & Your cross-project preferences \\ \hline
\texttt{\textasciitilde/.claude/skills/} & Global & Skills available in every project \\ \hline
\texttt{\textasciitilde/.claude/projects/\ldots/memory/} & Global & Auto memory: the tool's own per-project notes \\ \hline
\hline
\texttt{<project>/CLAUDE.md} & Local & Project brain, read every session \\ \hline
\texttt{<project>/.claude/settings.json} & Local & Project permissions (overrides global) \\ \hline
\texttt{<project>/.claude/skills/} & Local & This project's slash commands \\ \hline
\texttt{<project>/.claude/agents/} & Local & Specialized reviewer roles \\ \hline
\texttt{<project>/.claude/rules/} & Local & Always-on guardrails \\ \hline
\end{tabular}
\end{center}

\vspace{0.15cm}

\footnotesize\textbf{Other agents follow analogous layouts.} Codex CLI uses \texttt{AGENTS.md}; Gemini CLI uses \texttt{GEMINI.md}. The pattern --- global defaults plus a project-specific system prompt --- is durable across products.

\vspace{0.05cm}

\textbf{Rule:} global dotfolder $=$ your personal defaults (apply everywhere). Local dotfolder $=$ the project's rules (apply only here). T4 fills the local layer with research content.
\end{frame}

\begin{frame}{Initialization: What Happens When You Type \texttt{claude}}

Five steps run \textbf{before you type your first message} --- each one loads a file from the previous frame:

\vspace{0.2cm}

\begin{enumerate}
    \item \textbf{Read global config} --- \texttt{\textasciitilde/.claude/settings.json}: permission defaults, model preference

    \vspace{0.1cm}

    \item \textbf{Detect project context} --- searches current directory (and parents) for \texttt{CLAUDE.md} and \texttt{.claude/}

    \vspace{0.1cm}

    \item \textbf{Load project memory} --- reads \texttt{CLAUDE.md} (and the tool's auto memory for this project) into the context window

    \vspace{0.1cm}

    \item \textbf{Auto-load all rules} --- every \texttt{.claude/rules/*.md} file is loaded silently

    \vspace{0.1cm}

    \item \textbf{Ready} --- context window pre-populated before you say anything
\end{enumerate}

\vspace{0.2cm}

\textbf{Analogy:} steps 3--4 are a professor re-reading lecture notes before class. Students arrive; context is already loaded.

\vspace{0.1cm}

\textbf{Cost warning:} steps 3--4 consume tokens before you type anything. Keep \texttt{CLAUDE.md} and rules files concise.
\end{frame}

\begin{frame}{The CLI at a Glance: One Screen, Four Zones}
\textbf{Everything you will ever look at in a terminal-agent session is one screen with four zones.}

\vspace{0.1cm}

\begin{center}
\begin{tikzpicture}[
    zone/.style={rectangle, draw=black!50, align=left, font=\ttfamily\tiny, inner sep=3pt},
    lab/.style={font=\scriptsize, text=RedTitle, align=left},
    larr/.style={-{Stealth[length=1.6mm]}, thin, RedTitle!70}
]
    % terminal window
    \node[zone, fill=white, minimum width=8.6cm, minimum height=2.15cm, text width=8.3cm] (tr) at (0,0)
        {You> Parse the roster page into a seed CSV\\[1pt]
         claude> {[Read sources/fellows\_page.html]}\\
         \phantom{claude>} {[Write data/fellows\_seed.csv]}\\
         \phantom{claude>} Parsed 882 rows; flagged 2 unparseable\\
         \phantom{claude>} election years in notes/batch\_log.md};
    \node[zone, fill=orange!10, draw=orange!70, minimum width=8.6cm, minimum height=0.55cm, text width=8.3cm, below=0.06cm of tr] (perm)
        {Allow Bash(python scripts/parse\_roster.py)?\ \ [y]es / [n]o / always};
    \node[zone, fill=blue!4, minimum width=8.6cm, minimum height=0.5cm, text width=8.3cm, below=0.06cm of perm] (inp)
        {> \_};
    \node[zone, fill=VeryLightGray, minimum width=8.6cm, minimum height=0.45cm, text width=8.3cm, below=0.06cm of inp] (stat)
        {mode: plan\ \ $\cdot$\ \ context used: 34\%\ \ $\cdot$\ \ session cost so far: /cost};
    % labels
    \node[lab, right=0.45cm of tr.east, text width=3.2cm] (l1)
        {\textbf{1 Transcript}\\ \scriptsize dialogue + bracketed tool traces (the JSON calls, printed)};
    \node[lab, right=0.45cm of perm.east, text width=3.2cm] (l2)
        {\textbf{2 Permission prompt}\\ \scriptsize approve actions one by one};
    \node[lab, right=0.45cm of inp.east, text width=3.2cm] (l3)
        {\textbf{3 Your input}\\ \scriptsize questions, commands, /slash};
    \node[lab, right=0.45cm of stat.east, text width=3.2cm] (l4)
        {\textbf{4 Status line}\\ \scriptsize mode, context usage, cost};
    \draw[larr] (l1.west) -- (tr.east);
    \draw[larr] (l2.west) -- (perm.east);
    \draw[larr] (l3.west) -- (inp.east);
    \draw[larr] (l4.west) -- (stat.east);
\end{tikzpicture}
\end{center}

\vspace{0.1cm}
\footnotesize\textit{Zone 1's \texttt{[Read ...]} traces are the JSON tool calls of the earlier frame, printed for you --- the audit trail is on screen by default. Layout details differ across products and versions; the four zones are durable.}
\end{frame}

\begin{frame}{Commands, Modes, Keys: The Control Surface}

\begin{center}
\footnotesize
\begin{tabular}{|l|p{5.0cm}|p{3.3cm}|}
\hline
\textbf{Control} & \textbf{What it does} & \textbf{When to use} \\ \hline
\multicolumn{3}{|l|}{\textbf{Modes} (cycle with \texttt{Shift+Tab})} \\ \hline
Plan mode & Read-only: the agent explores and drafts a plan, executes \textbf{nothing} until you approve & Start every non-trivial task here \\ \hline
Accept-edits & Auto-approves routine file edits & Only after trust is earned \\ \hline
\multicolumn{3}{|l|}{\textbf{Slash commands}} \\ \hline
\texttt{/effort max} & Maximum reasoning depth: best quality, slower, costlier & Novel problems, economic correctness \\ \hline
\texttt{/effort default} & Balanced quality and speed & Routine work \\ \hline
\texttt{/cost} & Tokens used + estimated USD this session & Budget check before long tasks \\ \hline
\texttt{/compact}, \texttt{/clear} & Manage the context window & Next section of this deck \\ \hline
\multicolumn{3}{|l|}{\textbf{Keys}} \\ \hline
\texttt{Esc} & Interrupt the agent mid-task immediately & When it's going the wrong way \\ \hline
\end{tabular}
\end{center}

\vspace{0.05cm}

\textbf{The effort--cost tradeoff:} maximum-effort modes plus long context can be an order of magnitude more expensive than default. Spend effort on \textit{decisions}; save it on \textit{execution}.
\end{frame}

\begin{frame}[fragile]{Permissions: The Blast-Radius Dial}

\begin{columns}[T]
\begin{column}{0.44\textwidth}
\textbf{The agent's toolbox:}

\vspace{0.15cm}
\footnotesize
\begin{tabular}{|l|p{2.8cm}|}
\hline
\textbf{Tool} & \textbf{What it does} \\ \hline
\texttt{Bash} & Run any shell command \\ \hline
\texttt{Read}/\texttt{Grep} & Read and search files \\ \hline
\texttt{Write} & Create or overwrite a file \\ \hline
\texttt{Agent} & Spawn a child instance \\ \hline
\end{tabular}

\vspace{0.15cm}
\textbf{``Which Python ran this?''}
\begin{itemize}
    \item \scriptsize the agent's \texttt{Bash} inherits your shell's \texttt{\$PATH}: if you ran \texttt{conda activate myenv} before \texttt{claude}, \textit{that} Python runs
    \item \scriptsize add a \texttt{which python} check to \texttt{CLAUDE.md}
\end{itemize}
\end{column}
\begin{column}{0.53\textwidth}
\textbf{The permission model:}

\vspace{0.1cm}
\footnotesize
\begin{itemize}
    \item \textbf{Default: ask.} You see the exact command before it runs
    \item \textbf{Allowlist} trusted patterns in \texttt{.claude/settings.json}:
\end{itemize}
\begin{lstlisting}[style=pythonstyle,language={},basicstyle=\ttfamily\tiny,numbers=none]
"permissions": {
  "allow": ["Bash(python *)",
            "Bash(git *)"],
  "deny":  ["Read(.env)",
            "Bash(rm -rf *)"] }
\end{lstlisting}
\begin{itemize}
    \item \textbf{Bypass-permissions} skips all prompts: isolated sandboxes only
\end{itemize}
\end{column}
\end{columns}

\vspace{0.08cm}
\begin{conceptbox}
\footnotesize
\textbf{Think of authorizing a new RA.} Default $=$ every action needs your signature; allowlist $=$ standing approval for routine work; deny $=$ never allowed; bypass $=$ handing over your own login. \emph{``Blast radius''} $=$ how much damage one wrong action could do before you notice.
\end{conceptbox}
\end{frame}

%===========================================================
%=============================================================================
\section{Your First Session}
%===========================================================

\sectiondivider{Your First Session}{A real walkthrough on the fellows roster}

\begin{frame}[fragile]{First Session (1): Launch on a Real Project}
\textbf{Initialization in the abstract $\to$ what it actually looks like, on the fellows project from T1.}

\vspace{0.08cm}
\begin{lstlisting}[basicstyle=\ttfamily\scriptsize,numbers=none]
$ cd ~/research/es-fellows
$ claude
[reads ~/.claude/settings.json]
[reads ./CLAUDE.md]
[loads .claude/rules/provenance.md]
[ready]

You> Parse sources/fellows_page.html into data/fellows_seed.csv
     with columns: name, affiliation, election_year, profile_url.

claude> [Read sources/fellows_page.html]
        [Write data/fellows_seed.csv]
        Parsed 882 rows. Two election years did not parse
        (flagged in notes/batch_log.md). Spot-check 5 rows?
\end{lstlisting}

\vspace{0.05cm}
\footnotesize
\begin{enumerate}
    \item Initialization is automatic; no flags needed for the first session
    \item The agent reads context (rules, references) \emph{before} acting --- visible as bracketed \texttt{[Read ...]} traces
    \item Each trace is a JSON tool call that the harness translated into a real disk operation
    \item Note the agent's last line: a competent session \emph{ends turns with a checkable claim}, not just ``done''
\end{enumerate}
\end{frame}

\begin{frame}{First Session (2): Pen and Paper, Then Plan Mode}

\begin{columns}[T]
\begin{column}{0.52\textwidth}
\small
\begin{enumerate}
    \item[\textbf{0.}] \textbf{Before the agent---pen~+~paper}\\
    Write one sentence: \textit{``I want to compute X using method Y for question Z.''}\\
    {\footnotesize Fellows version: ``a provenance-aware seed roster of all current ES Fellows, from the official page only.''}

    \vspace{0.1cm}

    \item[\textbf{1.}] \textbf{Enter plan mode (\texttt{Shift+Tab})}\\
    Describe your goal. The agent drafts a plan. \textit{Nothing runs yet.}

    \vspace{0.1cm}

    \item[\textbf{2.}] \textbf{Review and approve}\\
    Read the plan. Ask questions. Revise. Then approve it.
\end{enumerate}
\end{column}
\begin{column}{0.44\textwidth}
\begin{tcolorbox}[colback=blue!3, colframe=RedTitle!60]
\footnotesize
\textbf{Why this matters:} plan review is the cheapest place to catch a wrong task definition.
\end{tcolorbox}
\vspace{0.2cm}
\footnotesize
\textbf{Common beginner mistake:} skipping Step~0. Without a clear goal, the agent will confidently build the wrong thing.

\vspace{0.15cm}
\textit{What exactly to check before approving a plan: T3 gives the five criteria.}
\end{column}
\end{columns}
\end{frame}

\begin{frame}{First Session (3): Implement, Review, Iterate}
\begin{columns}[T]
\begin{column}{0.52\textwidth}
\small
\begin{enumerate}
    \item[\textbf{3.}] \textbf{Implement}\\
    The agent executes the approved plan. Press \texttt{Esc} to stop if wrong.

    \vspace{0.1cm}

    \item[\textbf{4.}] \textbf{Review output}\\
    Does it make economic sense? Re-derive one number by hand.\\
    {\footnotesize Fellows version: open 5 random rows; follow each source URL.}

    \vspace{0.1cm}

    \item[\textbf{5.}] \textbf{Iterate}\\
    Fix, extend one feature, compact context, and return to step~1.
\end{enumerate}
\end{column}
\begin{column}{0.44\textwidth}
\begin{center}
\begin{tikzpicture}[
    stagebox/.style={rectangle, draw=black, fill=VeryLightGray, rounded corners,
                 minimum width=2.8cm, minimum height=0.65cm,
                 text width=2.6cm, align=center, font=\footnotesize},
    arrow/.style={-{Stealth[length=2mm]}, thick}
]
    \node[stagebox] (s0) {0.\ Pen + paper};
    \node[stagebox, below=0.22cm of s0] (s1) {1.\ Plan mode};
    \node[stagebox, below=0.22cm of s1] (s2) {2.\ Approve};
    \node[stagebox, below=0.22cm of s2] (s3) {3.\ Implement};
    \node[stagebox, below=0.22cm of s3] (s4) {4.\ Review};
    \draw[arrow] (s0) -- (s1);
    \draw[arrow] (s1) -- (s2);
    \draw[arrow] (s2) -- (s3);
    \draw[arrow] (s3) -- (s4);
    \draw[arrow] (s4.east) -- ++(0.6,0) -- ++(0,2.95) -- (s1.east);
    \node[right=0.4cm of s3, font=\tiny, rotate=90] {iterate};
\end{tikzpicture}
\end{center}
\vspace{0.2cm}
\footnotesize
\textbf{Practical takeaway:} after each implementation pass, force one human check before you ask for the next extension.
\end{column}
\end{columns}
\end{frame}

\begin{frame}{What the First Session Leaves Behind}
\textbf{Close the terminal. What survives is not the conversation --- it is the project state on disk.}

\vspace{0.2cm}

\begin{columns}[T]
\begin{column}{0.5\textwidth}
\textbf{In the folder now:}
\footnotesize
\begin{itemize}
    \item \texttt{data/fellows\_seed.csv} --- the artifact (882 rows)
    \item \texttt{notes/batch\_log.md} --- what parsed, what was flagged
    \item \texttt{CLAUDE.md} --- two new lines: the parse rules you settled on
    \item \texttt{.claude/rules/provenance.md} --- the source rule, now always-on
\end{itemize}
\end{column}
\begin{column}{0.46\textwidth}
\textbf{Why this is the whole game:}
\footnotesize
\begin{itemize}
    \item chat history is session-bound (resumable, but do not rely on it); \textbf{files persist and compose}
    \item next session, initialization reloads all of this \emph{automatically} --- batch 2 starts where batch 1 ended
    \item ``the project files hold the rules and state across sessions'' --- the public-web analog of \texttt{CLAUDE.md} (Lec10 T4)
\end{itemize}
\end{column}
\end{columns}

\vspace{0.25cm}

\begin{takeawaybox}
\footnotesize
\textbf{A session is good when it improves the project state, not when the conversation felt smart.}
\end{takeawaybox}
\end{frame}

\begin{frame}{Close the Loop: Reflect, Save, Shrink}
\textbf{The last five minutes of a session are worth more than the next fifty. Three habits.}

\vspace{0.2cm}

\begin{enumerate}
    \item \textbf{Reflect.} End with a two-minute review: what worked, what you corrected, what surprised you. Ask the agent itself to summarize the session's decisions --- it is cheap and surprisingly useful
    \medskip
    \item \textbf{Save.} Write that summary into the plan or progress file, and update \texttt{CLAUDE.md} while it is fresh. The next session then starts from the updated plan, not from anyone's memory
    \medskip
    \item \textbf{Shrink.} Prefer several small, focused sessions over one giant thread: a small context is cheaper, and the model's recall is sharper (mechanics in the next section). Break the big project into bounded sessions --- \emph{one intermediate product per session}
\end{enumerate}

\vspace{0.2cm}

\begin{takeawaybox}
\footnotesize
\textbf{A big project is a sequence of small sessions, each ending with an updated plan.} T4 turns this rhythm into a management system.
\end{takeawaybox}
\end{frame}

%===========================================================
%=============================================================================
\section{Budget and Discipline}
%===========================================================

\sectiondivider{Budget and Discipline}{Tokens, cost, the context window --- and the pillars that keep you honest}

\begin{frame}{Tokens and Cost on the Radar}
\textbf{Every agent turn is a model API call. Cost scales with input and output tokens.}

\vspace{0.2cm}

\begin{itemize}
    \item \textbf{Input tokens.} Your prompt, conversation history, \texttt{CLAUDE.md}, rules, retrieved chunks. Every step that adds context adds cost.
    \item \textbf{Output tokens.} Generated code, JSON tool calls, prose explanations. Reasoning-heavy modes also produce many internal tokens.
    \item \textbf{Long sessions accumulate context.} Watch the status line; check \texttt{/cost} before long tasks.
    \item \textbf{Pin for reproducibility.} Record the model version, prompt template, and seed when applicable; AEA disclosure (Oct 2024) extends to AI-assisted analysis.
\end{itemize}

\end{frame}

\begin{frame}{The Context Window Is the Agent's Working Memory}
\textbf{Everything the model ``knows'' during a turn sits in one finite window (Lec07 T3b). Here is what fills it in an agent session:}

\vspace{0.15cm}

\begin{center}
\begin{tikzpicture}[
    seg/.style={draw=black!55, minimum height=0.72cm, font=\tiny, align=center, inner sep=1.5pt},
    wcap/.style={font=\tiny\itshape, text=black!65, align=center}
]
    \node[seg, fill=blue!12,  minimum width=1.9cm, text width=1.8cm] (s1) at (0,0) {system prompt + tool schemas};
    \node[seg, fill=blue!7,   minimum width=2.1cm, text width=2.0cm, right=0cm of s1] (s2) {\texttt{CLAUDE.md} + rules + auto memory};
    \node[seg, fill=blue!4,   minimum width=1.3cm, text width=1.2cm, right=0cm of s2] (s3) {skill index};
    \node[seg, fill=green!8,  minimum width=2.6cm, text width=2.5cm, right=0cm of s3] (s4) {conversation so far};
    \node[seg, fill=orange!12, minimum width=3.6cm, text width=3.5cm, right=0cm of s4] (s5) {tool outputs: file contents, command results, errors};
    \node[wcap] at ($(s1.south)!0.5!(s3.south)+(0,-0.32)$) {loaded at initialization (the frame before)};
    \node[wcap] at ($(s4.south)!0.5!(s5.south)+(0,-0.32)$) {grows every turn $\longrightarrow$ dominates long sessions};
\end{tikzpicture}
\end{center}

\vspace{0.2cm}

\begin{itemize}
    \item The window is a \textbf{finite budget}: model-dependent, but always exhaustible by a long session
    \item Recall is not uniform --- material buried mid-window is recalled worst (``lost in the middle,'' Lec07 T3b); we do not reteach it, we \emph{manage around} it
    \item Consequence: \textbf{long sessions degrade} --- the agent slowly forgets early decisions and corrections
\end{itemize}
\end{frame}

\begin{frame}{Compaction: What \texttt{/compact} Keeps and What It Drops}
\textbf{When the window fills, the harness summarizes: that is compaction. It is a lossy operation, and you should treat it as one.}

\vspace{0.1cm}

\begin{itemize}
    \item \texttt{/compact} = \textbf{summarize and clear}: older tool outputs are dropped first, the dialogue is condensed into a summary, and \texttt{CLAUDE.md} + rules are re-read from disk
    \item \textbf{What survives:} the decisions, the current task, key snippets. \textbf{What often does not:} the \emph{reasoning} behind decisions, rejected alternatives, small corrections you made in passing
    \item \textbf{Auto-compaction:} by default the harness compacts on its own as the window fills --- convenient, but the same information loss applies whether you triggered it or not
    \item \texttt{/clear} = \textbf{full reset}: conversation gone, files and memory remain; right choice when switching to an unrelated task
\end{itemize}

\vspace{0.1cm}

\begin{takeawaybox}
\footnotesize
\textbf{Rule: state belongs in files, not in the conversation.} If a decision matters tomorrow, write it to \texttt{CLAUDE.md} or a progress note \emph{before} the window fills --- then compaction costs you nothing.
\end{takeawaybox}

\vspace{0.05cm}
\footnotesize\textit{That is the mechanics. T5 returns to what this trade-off costs research integrity --- and why more context is not always better.}
\end{frame}

\begin{frame}{The 5-Pillar Spine, One More Time}
\textbf{The five pillars (introduced in Lec01, exercised throughout Lec06) reappear here as the agent-era research workflow.}

\vspace{0.15cm}

\begin{columns}[T]
\begin{column}{0.34\textwidth}
\footnotesize
\begin{enumerate}
    \item \textbf{Theory}
    \item \textbf{Algorithm}
    \item \textbf{Numerics}
    \item \textbf{Advanced}
    \item \textbf{Code literacy}
\end{enumerate}
\vspace{0.1cm}
\scriptsize\textit{The agent accelerates each pillar; it replaces none of them.}
\end{column}
\begin{column}{0.63\textwidth}
\footnotesize
\begin{tabular}{|p{2.1cm}|p{5.7cm}|}
\hline
\textbf{Rendering} & \textbf{One-line discipline} \\ \hline
1. Question design & Economic language; define the evidence; limiting cases first \\ \hline
2. Procedure design & The right procedure \emph{and} its acceptance criterion, before code \\ \hline
3. Data/numerical discipline & Grids, weights, crosswalks, missingness --- saved, not implicit \\ \hline
4. Domain constraints & Know when the agent's convenient default is wrong \\ \hline
5. Code literacy & Read what was built; verify it is what you meant \\ \hline
\end{tabular}
\end{column}
\end{columns}

\vspace{0.1cm}
\textbf{The throughline:} pillars 1--4 decide \emph{what should be built and what counts as success}; the machine only decides how fast. Weak pillars $=$ faster nonsense.
\end{frame}

\begin{frame}{Division of Labor: You Provide, the Agent Provides}
\begin{center}
\small
\begin{tabular}{|p{5.5cm}|p{5.9cm}|}
\hline
\textbf{You provide} & \textbf{Agent provides} \\ \hline
Question, identification logic, and source rules & Fast code generation and refactoring \\ \hline
Workflow choice and benchmark design & Boilerplate implementation and debugging \\ \hline
Knowledge of the data, model, and institutions & File edits, script writing, and test runs \\ \hline
Economic interpretation and seminar defense & Iteration on visible errors and outputs \\ \hline
\end{tabular}
\end{center}

\vspace{0.25cm}

\textbf{Research version of the slogan:} the agent can be a strong RA, but you still have to be the principal investigator.

\end{frame}

\begin{frame}{Validation Is Your Referee Checklist}
\textbf{Every iteration needs both workflow checks and economic checks.}

\vspace{0.25cm}

\begin{columns}[T]
\begin{column}{0.48\textwidth}
\textbf{Workflow checks}
\begin{itemize}
    \item Convergence achieved
    \item No NaN or Inf values
    \item Mappings, weights, or source logs saved where expected
    \item Diagnostics, residuals, or quote checks pass
    \item Script reproduces the same output twice
\end{itemize}
\end{column}
\begin{column}{0.48\textwidth}
\textbf{Economic checks}
\begin{itemize}
    \item Limiting cases recover known results
    \item Comparative statics, subgroup patterns, or review claims make sense
    \item Magnitudes and source coverage are plausible
    \item Inference uses the right standard errors, sample, and restrictions
    \item You can explain the result in plain economic language
\end{itemize}
\end{column}
\end{columns}

\vspace{0.25cm}

\textbf{If you cannot defend the output in a seminar, you are not done.}
\end{frame}

\begin{frame}{Bridge: The Machine Needs a Method}
\textbf{You can now install the tool, launch it, steer it, budget it, and read its screen. What you cannot yet do is structure a week of research with it.}

\vspace{0.25cm}

\begin{itemize}
    \item Act 2 covered the \emph{mechanism}: components, tool calls, install, initialization, the CLI, permissions, the context window
    \item The next question: \emph{``how do I structure real work?''} --- what to ask for first, how big a step to take, when to trust a result
    \item Act 3's answer is the \textbf{homotopy workflow}: build the easy version, validate, then extend one margin at a time --- the exact discipline behind HA-Ramsey
\end{itemize}

\vspace{0.3cm}
\footnotesize\textit{Arc: T1 Framing $\to$ \textbf{T2 Under the Hood} $\to$ T3 Homotopy $\to$ T4 Project Org $\to$ T5 Mistakes/Context $\to$ T6 Wrap.}
\end{frame}

%===========================================================
\end{document}
